\section{Pruned encodings}
\label{sec:pruned}

The second summand of \eqref{eq:cost8} pays for the encodings of the input at
all \(10^{L}\) leaves.  Most of them are never read.  This section computes
only those that are.  First the mathematics: the correctness of the pruned
recursion, its number of operations, and the resulting cost of the method
(Theorem~\ref{thm:pruned}, Proposition~\ref{prop:costP}).  Then the program:
the pruned encoder is written in the language of \cite{upstream}, the
procedures of the paper's data structure are re-verified under the weaker
contract that the encodings are right only at the leaves that are read, and
the running time of the new solver is proved on the machine
(Theorem~\ref{thm:encoder}).  Everything is machine-checked.

\subsection{What is read and what is nonzero}

Let \(a\) be the array of a band.  It is indexed by the \emph{left strings}:
strings \(u=u_1\cdots u_L\) of left variables.  For a leaf
\(\tau=\tau_1\cdots\tau_L\) the encoding is
\[
  \Phi_\tau(a)=\sum_u a[u]\prod_{\ell=1}^{L}\varphi_{\tau_\ell}(u_\ell),
\]
where \(\varphi_\tau\) is the linear form of the term \(\tau\) in the left
variables.  Each \(\varphi_\tau\) has at most three nonzero coefficients.  Two
facts hold in the recursion.
\begin{itemize}
\item Only leaves with at most \(m\) symbols \(\Pz\) are read: the leaves of
  \(\Leaves(U)\) have order at least \(0\).
\item \(a[u]=0\) unless \(u\) has exactly \(m\) inner variables.
\end{itemize}

Yates' algorithm computes the encodings level by level.  Stage \(k\) holds
\[
  A_k[\tau_1\cdots\tau_k;u_{k+1}\cdots u_L]
  =\sum_{u_1,\dots,u_k}a[u]\prod_{\ell\le k}\varphi_{\tau_\ell}(u_\ell),
\]
so \(A_0=a\) and \(A_L[\tau]=\Phi_\tau(a)\), and
\(A_{k+1}[\tau_1\cdots\tau_{k+1};\cdot]
=\sum_v\varphi_{\tau_{k+1}}(v)A_k[\tau_1\cdots\tau_k;v\,\cdot]\).

\begin{definition}\label{def:pruned}
Let \(P_k\) be the set of strings of \(k\) terms with at most \(m\) symbols
\(\Pz\), and \(U_k\) the set of strings of \(L-k\) left variables with at most
\(m\) inner variables.  The \emph{pruned recursion} computes at stage \(k\)
only the entries of \(A_k\) with index in \(P_k\times U_k\), by the rule above,
and treats all other entries as \(0\).
\end{definition}

\subsection{Correctness and cost}

\begin{theorem}[Pruned encodings]\label{thm:pruned}
Let \(L\ge10m\), and let \(a\) vanish off the left strings with exactly \(m\)
inner variables.
\begin{enumerate}
\item The pruned recursion computes \(\Phi_\tau(a)\) for every leaf \(\tau\)
  with at most \(m\) symbols \(\Pz\).
\item \(|P_k|=\sum_{j\le m}\binomial kj9^{k-j}\) and
  \(|U_k|=\sum_{i\le m}\binomial{L-k}{i}4^{i}3^{L-k-i}\).
\item \(\sum_{k=0}^{L}|P_k|\,|U_k|\le\frac92|P_L|\le\frac92\cdot\frac{M}{1-\rho}\)
  with \(\rho=9m/(L-m+1)\).
\item Every computed entry is a sum of at most three products.  So the pruned
  recursion makes at most \(\frac{27}{2}M/(1-\rho)\) multiplications.
\end{enumerate}
\end{theorem}

\begin{proof}
(i) An entry of \(A_k\) whose suffix has more than \(m\) inner variables is
\(0\): every \(u\) in its sum has more than \(m\) inner variables, so
\(a[u]=0\).  An entry of \(A_{k+1}\) with prefix in \(P_{k+1}\) reads entries of
\(A_k\) whose prefix is in \(P_k\), since deleting a term does not add a
symbol \(\Pz\).  So, by induction on \(k\), the pruned arrays agree with those
of Yates' algorithm at every index whose prefix has at most \(m\) symbols
\(\Pz\): such an entry is \(0\) if its suffix is outside \(U_k\), and otherwise
each entry it reads either lies in \(P_{k-1}\times U_{k-1}\) or is \(0\) in both
arrays.

(ii) There are nine terms other than \(\Pz\), and four inner and three outer
left variables.

(iii) Since \(\binomial{k+1}{j}\ge\binomial kj\), we have
\(|P_{k+1}|\ge9|P_k|\).  By Pascal's rule,
\[
  \sum_{i\le m}\binomial{n+1}{i}4^{i}3^{n+1-i}
  =3\sum_{i\le m}\binomial ni4^{i}3^{n-i}
   +4\sum_{i\le m-1}\binomial ni4^{i}3^{n-i},
\]
so \(|U_k|\le7|U_{k+1}|\).  With \(|U_L|=1\) this gives
\(|P_k||U_k|\le(7/9)^{L-k}|P_L|\), and \(\sum_{k}(7/9)^{L-k}\le9/2\).  The terms
of \(|P_L|=\sum_{j\le m}\binomial Lj9^{L-j}\) have consecutive ratios
\(\binomial L{j-1}9/\binomial Lj=9j/(L-j+1)\le\rho<1\), and the term with
\(j=m\) is \(M\).  So \(|P_L|\le M/(1-\rho)\).

(iv) follows from (iii) and the three coefficients of a form.
\end{proof}

In Lean (directory \lean{PrunedEncoding}): Yates' algorithm is
\lean{yatesStage} with \lean{yatesStage_final}; the pruned recursion is
\lean{prunedStage}; part~(i) is \lean{prunedStage_eq_yatesStage} and
\lean{prunedStage_final}; part~(ii) is \lean{card_stageSet_of_le} with
\lean{prefCount} and \lean{sufCount}; part~(iii) is
\lean{sum_card_stageSet_le}; part~(iv) is \lean{prunedMults_le_M}; the theorem
as a whole is \lean{pruned_encoding}.

\begin{remark}\label{rem:PL}
\(|P_L|=\sum_{d=0}^{m}\beta_d\): the pruned recursion computes exactly the
encodings at the leaves of order at least \(0\).  For \(L\ge cm\) with \(c>10\)
fixed, \(1/(1-\rho)\le(c-1)/(c-10)\).
\end{remark}

\begin{remark}[The work of the program]\label{rem:encwork}
The program of Section~\ref{sec:encoder-program} below follows the recursion
of Theorem~\ref{thm:pruned} with one simplification: at stage \(k\) it keeps
the entries with index in \(P_k\times(\text{all \(7^{L-k}\) suffixes})\)
instead of \(P_k\times U_k\), so that the suffixes are indexed densely and no
compact indexing of \(U_k\) is needed.  For \(L\ge10m\) the largest slice,
with \(7^{L}\) entries, is no larger than a tile: \(7^{L}\le9^{L-m}\le M\)
(\lean{seven_pow_le_M}).  The number of computed entries, \lean{encWorkP}~\(L\,m\)
in Lean, is
\[
  \sum_{k=0}^{L}|P_k|\,7^{L-k}\le\frac92|P_L|
  \le\frac92\cdot\frac{M}{1-\rho},
\]
from \(9|P_k|\le|P_{k+1}|\) and \(\sum_i(7/9)^i\le9/2\)
(\lean{encWorkP_le}, \lean{encWorkP_le_M_div}): part~(iii) with \(7^{L-k}\)
in place of \(|U_k|\), and the same bound.  The time of the program is at most
\(\EncConst\) times this number (\lean{tEncP_le}).
\end{remark}

\subsection{The thinness bound}

With \(O(L)\) word operations per entry for the index arithmetic, a band costs
\(O(LM)\) instead of \(O(L\,10^{L})\), and all \(N/(\sqrt K N_0)\) bands cost
\(O(L\,N\sqrt K N_0)\), since \(M=KN_0^{2}\).  This is at most
\(N^{2}/D^{\gamma}\) up to the factor \(L\) when \(\sqrt K N_0D^{\gamma}\le N\).
So the thinness condition \eqref{eq:Rc} becomes
\begin{equation}\label{eq:RcP}
  \varepsilon<\RcP(\gamma):=\frac{\ln 4}{\Apr(c)+\gamma\ln 4},\qquad
  \Apr(c)=\frac c2\Hent\Bigl(\frac1c\Bigr)+(c-1)\ln 3 ,
\end{equation}
where \(\Apr(c)\) is the exponential rate of \(\sqrt K N_0\) per inner level.
Since \(M\le10^{L}\), we have \(\Apr(c)\le\Afull(c)\) and
\(\RcP(\gamma)\ge\Rc(\gamma)\).  In Lean the constants are \lean{basePruned} and
\lean{baseFull}, the bound is \lean{RcP}, and the two inequalities are
\lean{basePruned_le_baseFull} and \lean{Rc_le_RcP}.

\begin{proposition}[The costs with pruned encodings]\label{prop:costP}
Let \(c>10\), \(0<\theta<0.9\), \(0\le\gamma<\gX(c,\theta)\) and
\(0<\varepsilon<\RcP(\gamma)\).  There are \(C\) and \(m_0\) such that the
following holds for all \(D\ge2\) and \(N\) with \(D\le N^{\varepsilon}\) and
\(m=\lceil\log_4 D\rceil\ge m_0\), with \(L=\lceil cm\rceil\) and
\(t=\lceil\theta m\rceil\).
\begin{enumerate}
\item The tile fits: \(\sqrt K N_0\le N\).
\item The cost of the preprocessing with pruned encodings,
  \[
    L\,m\,\frac{\sum_{d=t}^{m}\beta_d}{M}\,N^{2}
    +\frac{N}{\sqrt K N_0}\cdot L\cdot\frac{27}{2}\cdot\frac{M}{1-\rho},
  \]
  is at most \(C\,N^{2}\log^{2}D/D^{\gamma}\).
\item The cost \(L\sum_{d\le t}\alpha_d\) of a query is at most
  \(C\,D^{q(\theta)}\log D\).
\end{enumerate}
\end{proposition}

\begin{proof}
The first summand of (ii) is bounded as in Section~\ref{sec:exact}.  The second
is \(\frac{27}{2}NL\sqrt K N_0/(1-\rho)\).  Now
\(\sqrt K N_0\le e^{m\Apr(L/m)}\), the ratio \(L/m\) tends to \(c\), and
\(N\ge D^{1/\varepsilon}>4^{(m-1)/\varepsilon}\); so
\(D^{\gamma}L\sqrt K N_0\le N\) for all large \(m\), because
\(\Apr(c)+\gamma\ln 4<\ln 4/\varepsilon\).  This gives (i) and (ii).  Part~(iii)
is as in the paper.
\end{proof}

In Lean (file \lean{PrunedEncoding/Costs.lean}): the expression of (ii) is
\lean{cost8P}, the inequality \(D^{\gamma}L\sqrt K N_0\le N\) is \lean{eq_10P},
and the proposition is \lean{costsP}.  It is the analogue of the statement
\lean{costsX} on which Theorem~\ref{thm:offline} rests, with the pruned
encoder's count of operations in place of \(10^{L}\) and \(\RcP\) in place of
\(\Rc\).

\begin{theorem}[The encoder program]\label{thm:encoder}
For all rational \(c=a/b>10\) and \(\theta=p/q\in(0,0.9)\), every
\(0<\gamma<\gX(c,\theta)\) and every \(\varepsilon<\RcP(\gamma)\), there is a
procedure, in the programming language of \cite{upstream}, that solves the thin
product on all instances and takes at most a constant times
\[
  \frac{N^{2}(\log D+1)^{2}}{D^{\gamma}}+w\,D^{q(\theta)}(\log D+1)
\]
steps on every instance with \(1\le D\le N^{\varepsilon}\) and at most \(w\)
wanted positions.
\end{theorem}

Theorem~\ref{thm:encoder} is the conclusion of Theorem~\ref{thm:offline} with
\(\RcP\) in place of \(\Rc\).  On paper it follows from
Theorem~\ref{thm:pruned} and Proposition~\ref{prop:costP}: replace the encoder
of the offline routine by the pruned recursion, and replace the regime test of
the program, which is tied to \(\varepsilon\le1/18\), by the test with a
rational exponent of Section~\ref{sec:regime}.  In Lean the statement is the
proposition \lean{PrunedEncoderClaim} (file \lean{Pipeline/General.lean}),
proved as \lean{prunedEncoderClaim} (file \lean{Pipeline/PrunedClaim.lean}).
It is stated with the same notion, \lean{ThinClaim}, in which
Theorems~\ref{thm:T3} and~\ref{thm:T4} obtain the thin product from the
existing program (\lean{thinClaim_gammaX}, \lean{thinClaimRS_gammaX}): a claim
about the time model
\lean{lightModel} of the upstream programs, which \cite{upstream} compiles to
the word RAM.  The rest of this section describes the program and its proof.

\subsection{The encoder program}
\label{sec:encoder-program}

\emph{The encoder.}  Upstream's encoder is a recursive procedure
\lean{encodeBody} that descends the ten terms of the identity level by level and
fills the \(10^{L}\) encodings of a band.  The pruned encoder
(\lean{encodePBody}, file \lean{PrunedProgram/Encode/Recursion.lean}) is the
same procedure with a budget \(j\) of symbols \(\Pz\): a loop over the nine
terms \(\lambda\ne\Pz\) recurses with the same budget, and the term \(\Pz\)
(digit~\(9\)) is treated after the loop, only if the budget is nonzero, with
budget \(j-1\).  The last term is treated outside the loop because the
language's \texttt{for} charges every round the same time.  The slices are the
dense ones of Remark~\ref{rem:encwork}.  Its correctness is \lean{encodeP_spec}:
after the call with budget \(j\), every leaf with at most \(j\) symbols \(\Pz\)
holds its encoding (the postcondition \lean{EncodePostP}), and its time is at
most \(\EncConst\) times the work \lean{encWorkP}~\(L\,j\) (\lean{tEncP_le}),
which Remark~\ref{rem:encwork} bounds by
\(\frac92\cdot\EncConst\cdot M/(1-\rho)\) for \(j=m\).

\emph{The weaker contract.}  Upstream's data structure assumes that the
encodings of a band are right at all \(10^{L}\) leaves (a \emph{segment}).
With the pruned encoder they are right only at the leaves with at most \(m\)
symbols \(\Pz\); this is the predicate \lean{SegOn} (file
\lean{PrunedProgram/SegOn.lean}).  The preprocessing and the query of
Theorem~30 of the paper read an encoding cell in exactly two places, both at
such leaves: in the construction of the trie, at the leaf of a code of
\lean{nineStrs}, which has at most \(m-t\) digits~\(9\); and in a query, at a
leaf that contributes to the output string, which has the symbol \(\Pz\) only
at inner levels (\lean{P0Levels_subset_innerSetO},
\lean{card_P0Levels_le_of_contributes}).  So upstream's own procedures for the
tiles (\lean{fillListBody}, \lean{tileBody}, \lean{allTilesBody}) and for a query
(\lean{queryAtBody}) are re-verified under the weaker contract, with their program
text unchanged: \lean{fillListP_meets} and \lean{preCoreP_of_base58} (files
\lean{PrunedProgram/Tiles}, \lean{PrunedProgram/Pre/Preprocessing.lean}),
\lean{queryCoreP_spec} and \lean{queryAtP_of_base58} (directory
\lean{PrunedProgram/Query}).  Upstream's shared stage, which builds the
tables and encodes every band, becomes \lean{sharedPBody} with the entry
\lean{sharedP_entry} (file \lean{PrunedProgram/Shared/Stage.lean}): its two
calls of the encoder are exchanged for calls of the pruned one, and nothing
else changes.  The invariant of the finished data structure is \lean{DSReadyP}
(file \lean{PrunedProgram/Pre/Defs.lean}).

\emph{The program.}  The routines with rational parameters are
\lean{pre31PBody}, \lean{offline32PBody} and \lean{offline32P_meets}
(directory \lean{PrunedProgram/Wrapper}).  The program is \lean{programP}: the
program \lean{programRS} of Section~\ref{sec:regime} followed by seven
procedures, \lean{encodeP}, \lean{encodeBandsP}, \lean{sharedP},
\lean{preCoreP}, \lean{pre31P}, \lean{offline32P} and \lean{allP}; every other
procedure is upstream's, with its number.  The offline routine as an object of
the pipeline is \lean{offline32PRoutine}.

\emph{The cost.}  The running time is analyzed against \lean{cost8P} as the
running time of upstream's program is analyzed against \lean{cost8X} in
Theorem~\ref{thm:offline}: \lean{Within8P} and \lean{exists_tPreCore_leP},
\lean{CostsWithinP} and \lean{offlineWithinP} (directory \lean{Cost8P}) are
the analogues of \lean{exists_tPreCore_leX}, \lean{costsX} and
\lean{offlineWithinX}.  The small summands of the shared stage are bounded by
\(NL\sqrt K N_0\) directly, since \lean{cost8P} has no term \(10^{L}\).

\emph{The pipeline.}  The solver of all instances is generic in the offline
routine: the structure \lean{OfflineRoutine} (file
\lean{Pipeline/ThinClaim.lean}) packages the procedure, its time, what it
needs from the program and its specification, as the regime test is already
packaged as \lean{RegimeTest}.  Upstream's routine is the instance
\lean{offline32Routine} and the pruned one is \lean{offline32PRoutine};
\lean{thinClaimRS_of_offlineWithinT} (file \lean{Pipeline/RegimeRS.lean})
turns either into a thin-product claim at the regime \(D^{r}\le N^{s}\).  For
\lean{thinClaimP_gammaX} and \lean{prunedEncoderClaim} the fraction \(r/s\) is
chosen strictly between \(\Apr(c)/\ln 4\) and \(1/\varepsilon\), which is
possible because \(\varepsilon<\RcP(\gamma)<\ln 4/\Apr(c)\); the case
\(\varepsilon\le0\) of the statement is reduced to a positive \(\varepsilon\)
by \lean{ThinClaim.mono}.
